% titel: Blatt 8, A5: Wasser-Phasendiagramm (Klausuraufgabe, wahr/falsch)
% skript: Mortimer Kap. 11 (Phasendiagramm Wasser)

% Teil I: Werkzeuge, kein Lösungsweg.

Diese Aufgabe prüft, ob du das Phasendiagramm von Wasser quantitativ lesen
kannst (Aussagen als wahr/falsch kreuzen).

\begin{werkzeugkasten}[Wasser im p,T-Diagramm]
Wichtige Zahlen für Wasser:
\begin{itemize}
  \item Normalsiedepunkt des Wassers: $100\,$°C bei $p = 1.013\ bar$
        ($= 101.3\ kPa$).
  \item \textbf{Tripelpunkt} Wasser: $T = 0.01\,$°C und $p = 0.0061\ bar
        = 0.61\ kPa$. Dort koexistieren Eis, Wasser, Wasserdampf.
  \item Die fest/flüssig-Linie des Wassers fällt nach links:
        der Gefrierpunkt \emph{sinkt} mit steigendem Druck (Anomalie, dichte
        Flüssigkeit).
  \item Unterhalb des Tripelpunkts ($p < 0.61\ kPa$) gibt es kein flüssiges
        Wasser --- Sublimation!
\end{itemize}
\textit{Skript:} Mortimer Kap. 11 (Tabelle der Phasendiagrammdaten).
\end{werkzeugkasten}

\begin{achtung}[Darauf achten]
\begin{itemize}
  \item Um eine Aussage zu beurteilen, vergleicht man die Zahlen mit den
        bekannten Normalbedingungen und dem Tripelpunkt.
  \item ``Sieden bei 100\,°C`` gilt nur bei $1\ atm$, nicht bei jedem Druck.
  \item Frage dich bei jeder Aussage: Welche Punkte der $p$-$T$-Fläche kenne
        ich exakt, und welche Steigung haben die Linien?
\end{itemize}
\end{achtung}
